\chapter{Full Proof of the Composition Failure Theorem}
\label{app:proof}

\begin{quotation}
\noindent\textit{Synopsis.} This appendix gives a complete proof of
Theorem~\ref{thm:composition-failure}, together with the explicit
construction on which it rests. The proof is built to satisfy the
robustness requirement raised in \cref{rem:discard-rules-diverge}: rather
than exhibiting a witness that happens to work under one convenient local
discard rule, \cref{lem:exhaustion} below shows that the witnessing
relay's key discard is forced under \emph{any} construction obeying
\cref{def:budget-discard}, closing that specific open question for the
purposes of this theorem even though the broader question of whether
every claim in this monograph enjoys the same robustness remains open
(\cref{app:open-questions}).
\end{quotation}

\section{Recovery Tasks and the Witnessing Task Family}

The theorem requires an explicit, non-enumerable-in-advance family of
tasks $T$, together with a specific $\tau^* \in T$ for which insufficiency
can be verified directly rather than merely asserted. Both are supplied by
a single, simple family of tasks built directly from the minimal-pair
machinery of Chapter~\ref{ch:distinctions}.

\begin{definition}[Recovery Task]
\label{def:recovery-task}
For a named distinction $d$ with minimal pair $\{x_d, x_{\neg d}\}$
(\cref{def:minimal-pair}), the recovery task $\tau_d$ is the task of
correctly reporting, given a substrate, whether the source object $e$
satisfies $x_d$ or $x_{\neg d}$. Its correctness set
(\cref{def:correctness-set}) is the singleton containing whichever of the
two reports matches $e$'s actual state.
\end{definition}

\begin{lemma}[Sufficiency for a Recovery Task Iff No Discard]
\label{lem:recovery-sufficiency}
A rendering $r$ is sufficient (\cref{def:sufficiency}) for $\tau_d$ if and
only if $d \notin \discard{r}$.
\end{lemma}

\begin{proof}
If $d \notin \discard{r}$, the pair $\{x_d, x_{\neg d}\}$ is not
identified by $r$ (\cref{def:discarded-pair}, \cref{def:discard-set}), so
$r$ distinguishes which of the two states $e$ occupies, and the correct
report is the unique continuation reachable: $\admis{\tau_d}{r} =
\mathcal{C}_{\tau_d}(e)$, satisfying \cref{def:sufficiency}. If $d \in
\discard{r}$, the pair is identified, so a consumer with access only to
$r$ cannot rule out either report: both the correct and the incorrect
report are reachable, $\admis{\tau_d}{r} \supsetneq \mathcal{C}_{\tau_d}(e)$,
and \cref{def:sufficiency}'s containment fails.
\end{proof}

\begin{definition}[Witnessing Task Family]
\label{def:witnessing-family}
$T_{\mathrm{rec}} := \{\tau_d : d \text{ ranges over the open-ended space
of possible named distinctions}\}$.
\end{definition}

$T_{\mathrm{rec}}$ is open in the sense required by
\cref{prop:task-unavailability}: the argument given there --- that a task
requiring a distinction not represented in a fixed candidate list can
always be constructed after the fact --- applies verbatim to
$T_{\mathrm{rec}}$, since the space of possible minimal pairs
(\cref{def:minimal-pair}) that might later be named as relevant is not
fixed in advance by any finite specification of a substrate chain's
producers.

\section{A Construction-Independent Exhaustion Lemma}

\cref{rem:discard-rules-diverge} showed that \cref{def:budget-discard} is
satisfied by more than one non-equivalent construction. The following
lemma isolates a sufficient condition for a discard to be forced under
every such construction, so that the theorem's witness does not depend on
which one a given stage happens to use.

\begin{lemma}[Exhaustion Forces Discard]
\label{lem:exhaustion}
Fix a stage with budget $\beta_i$ and suppose the distinctions required by
$\tau_i$ (in the sense of \cref{def:benign-loss}) have total cost
$R_i = \sum_{d \text{ required}} C_i(d) \le \beta_i$. If every
non-required distinction $d'$ satisfies $C_i(d') > \beta_i - R_i$, then
every non-required distinction is dropped -- $d' \in D_i$ for every such
$d'$ -- under \emph{any} construction satisfying \cref{def:budget-discard}
together with the requirement that required distinctions are ranked
strictly above non-required ones.
\end{lemma}

\begin{proof}
Any construction satisfying \cref{def:budget-discard} must respect the
budget constraint $\sum_{d \in S} C_i(d) \le \beta_i$ for the funded set
$S$, and, by hypothesis, must rank every required distinction above every
non-required one. Since the required distinctions have total cost
$R_i \le \beta_i$, funding all of them is feasible and, given their
priority, is favored by both \cref{con:greedy-discard} (which processes
higher-priority items first and encounters no obstruction, since $R_i \le
\beta_i$) and \cref{con:optimal-discard} (for which funding every required
item dominates any allocation that sacrifices one, provided required
items are assigned value strictly exceeding any combination of
non-required items, which the priority requirement above entails). After
funding the required set, the remaining budget is exactly $\beta_i - R_i$.
By hypothesis, $C_i(d') > \beta_i - R_i$ for every non-required $d'$, so
no non-required distinction can be funded in addition to the required set
under either construction, nor under any construction obeying the same
budget constraint, since funding any single such $d'$ alongside the
required set would require total expenditure exceeding $\beta_i$. Hence
$d' \in D_i$ for every non-required $d'$, regardless of which construction
produced the allocation.
\end{proof}

\cref{lem:exhaustion} is the appendix's answer to the robustness question
raised in Chapter~\ref{ch:budget-relay}: it does not show that every
budget relay behaves identically under every construction, but it
identifies a class of stages -- those where the required set exactly, or
near-exactly, exhausts the budget -- whose discard behavior is
construction-independent, and the witness built below is chosen to fall
into exactly this class.

\section{Recalling Local Admissibility}

Chapter~\ref{ch:composition-failure} makes precise, in
\cref{def:locally-admissible}, what Theorem~\ref{thm:composition-failure}
means by a transition being ``locally admissible with respect to its own
stage's task $\tau_i$.'' That definition does not privilege
\cref{con:greedy-discard} over \cref{con:optimal-discard}, or vice versa;
it constrains only the priority ordering's relationship to $\tau_i$'s own
stated requirements, consistent with \cref{rem:discard-rules-diverge}'s
observation that \cref{def:budget-discard} constrains outcomes, not
procedures. It is used throughout the remainder of this appendix exactly
as stated there.

\section{The Canonical Witness}

\begin{construction}[Two-Stage Witness]
\label{con:witness}
Let $A_0$ carry five distinctions $d_1, \ldots, d_5$ with costs
$C(d_1)=1$, $C(d_2)=2$, $C(d_3)=3$, $C(d_4)=2$, $C(d_5)=4$ -- the same
instance introduced in \cref{sec:worked-relay} and
\cref{sec:redistribution-worked}. Let $\tau_0$ require exactly $\{d_1,
d_2\}$, with $\beta_0 = 4$. Let $\tau_1$ require exactly $\{d_1, d_2\}$
as well, with $\beta_1 = 3$.
\end{construction}

\begin{proposition}[The Witness Is Locally Admissible and Forces the Same Discard Under Any Construction]
\label{prop:witness-properties}
In \cref{con:witness}: (a) $B_0$ and $B_1$ are locally admissible
(\cref{def:locally-admissible}) with respect to $\tau_0$ and $\tau_1$
respectively, under any construction satisfying \cref{def:budget-discard};
and (b) $D_0 = \{d_3, d_4, d_5\}$ and $D_1 = \emptyset$ regardless of
which such construction is used.
\end{proposition}

\begin{proof}
For stage $0$: the required set $\{d_1, d_2\}$ has total cost $R_0 =
1 + 2 = 3 \le \beta_0 = 4$. Every non-required distinction has cost
exceeding $\beta_0 - R_0 = 1$: $C(d_3)=3>1$, $C(d_4)=2>1$, $C(d_5)=4>1$.
\cref{lem:exhaustion} applies directly, giving $D_0 = \{d_3,d_4,d_5\}$
under any construction satisfying \cref{def:budget-discard} with required
distinctions ranked above non-required ones -- which is exactly
\cref{def:locally-admissible}'s condition, so (a) and (b) both hold for
stage $0$.

For stage $1$: $A_1$ retains only $\{d_1, d_2\}$ (nothing else survived
stage $0$), both required by $\tau_1$, with total cost $R_1 = 3 = \beta_1$
exactly. There are no non-required distinctions present at this stage to
consider, so the exhaustion hypothesis of \cref{lem:exhaustion} is
vacuously about an empty set, and any construction funds $\{d_1,d_2\}$ in
full (the only feasible allocation using the entire budget), giving
$D_1 = \emptyset$ regardless of construction. Local admissibility for
stage $1$ follows immediately, since there is nothing non-required
present to rank incorrectly.
\end{proof}

\section{Extension to Arbitrary Relay Length}

\begin{construction}[Pass-Through Stage]
\label{con:pass-through}
A stage $j$ is a pass-through stage if $\tau_j$ requires exactly the
distinctions present in $A_j$, and $\beta_j$ equals their total cost
exactly.
\end{construction}

\begin{proposition}[Extension to Arbitrary Length]
\label{prop:extension}
For every $n \ge 2$, there is an $n$-stage budget relay, every transition
of which is locally admissible, with $D(A_0 \Rightarrow A_n) =
\{d_3,d_4,d_5\}$ exactly as in \cref{con:witness}.
\end{proposition}

\begin{proof}
Take \cref{con:witness}'s two stages as stages $0$ and $1$, and insert
$n-2$ pass-through stages (\cref{con:pass-through}) at positions $2,
\ldots, n-1$, each requiring exactly $\{d_1,d_2\}$ at cost $3$, matching
$A_1$'s (and every subsequent stage's) content exactly. Each pass-through
stage is locally admissible by the same vacuous argument given for stage
$1$ in the proof of \cref{prop:witness-properties}, and discards nothing,
so $D_j = \emptyset$ for $j \ge 2$. By \cref{def:cumulative-discard},
$D(A_0 \Rightarrow A_n) = D_0 \cup D_1 \cup \cdots \cup D_{n-1} =
\{d_3,d_4,d_5\} \cup \emptyset \cup \cdots \cup \emptyset =
\{d_3,d_4,d_5\}$.
\end{proof}

\section{Proof of Theorem~\ref{thm:composition-failure}}

\begin{proof}
Fix any $n \ge 2$ and take the $n$-stage relay of
\cref{prop:extension}, with every transition locally admissible by
construction. Let $\tau^* := \tau_{d_5} \in T_{\mathrm{rec}}$
(\cref{def:recovery-task}, \cref{def:witnessing-family}). By
\cref{prop:witness-properties} and \cref{prop:extension}, $d_5 \in
D(A_0 \Rightarrow A_n)$, so $d_5 \in \discard{A_n}$ relative to $A_0$. By
\cref{lem:recovery-sufficiency}, $A_n$ is insufficient for $\tau_{d_5}$.

It remains to check $\tau^* \notin \{\tau_0, \ldots, \tau_{n-1}\}$. Every
$\tau_i$ in the construction is a preservation task, requiring exactly
$\{d_1,d_2\}$, with correctness set consisting of continuations that
retain both; $\tau^* = \tau_{d_5}$ is a recovery task in the sense of
\cref{def:recovery-task}, with correctness set consisting of a single
report about $d_5$ alone. These are distinct formal objects with distinct
correctness sets for any nontrivial choice of underlying object $e$, so
$\tau^* \neq \tau_i$ for any $i$.

$T_{\mathrm{rec}}$ is an open task space containing $\tau^*$, as argued
following \cref{def:witnessing-family}. This exhibits, for every $n \ge
2$, a budget relay and an open task space $T = T_{\mathrm{rec}}$ satisfying
every clause of Theorem~\ref{thm:composition-failure}.
\end{proof}

\section{Verifying the Corollaries}

\begin{proof}[Verification of \cref{cor:local-improvement}]
Suppose $\beta_1$ is increased in the witness of \cref{con:witness},
holding $\beta_0 = 4$ fixed. Stage $1$ receives only $A_1 = \{d_1,d_2\}$
regardless of $\beta_1$'s value, since $d_3, d_4, d_5$ were already
absent from $A_1$ by the time stage $1$ acts; no increase in $\beta_1$
can fund a distinction that is not present to be funded. $D(A_0
\Rightarrow A_n)$ is therefore unaffected by any change to $\beta_1$ (or
to any $\beta_j$ for $j \ge 1$), confirming that only an improvement to
$\beta_0$ specifically -- the stage where the loss actually occurred --
could have prevented the eventual insufficiency, and even that requires
$\tau_0$'s producer to have known, at stage $0$, that $d_5$ would later
matter: precisely the knowledge \cref{prop:task-unavailability} shows is
unavailable in general.
\end{proof}

\begin{proof}[Partial verification of \cref{cor:no-free-lunch}]
The existence half of \cref{cor:no-free-lunch} -- that cumulative discard
is non-decreasing in relay length -- is exactly
\cref{prop:monotonicity}, and \cref{prop:extension} shows this bound is
not merely non-strict but can be realized with equality across
arbitrarily many additional stages (the pass-through stages contribute
no further discard). This appendix does not, however, establish the
stronger, quantitative claim that cumulative discard \emph{risk} grows
across relays in general, only that it does not shrink in this
particular family of constructions. The tightness question --- whether a
sharper growth bound holds under specific coordination mechanisms --
remains open, exactly as flagged in Appendix~\ref{app:open-questions},
and is not resolved here.
\end{proof}

\section{What This Proof Does Not Establish}

This appendix proves Theorem~\ref{thm:composition-failure} and
\cref{cor:local-improvement} in full, and establishes the existence half
of \cref{cor:no-free-lunch}. It does not establish that every budget
relay of practical interest falls into the construction-independent class
identified by \cref{lem:exhaustion}; relays in which the required set
does not closely exhaust the local budget may genuinely behave
differently under \cref{con:greedy-discard} than under
\cref{con:optimal-discard}, and nothing here rules that out for such
cases. It also does not establish \cref{cor:no-free-lunch}'s quantitative
tightness, address the infinite-relay case flagged in
Chapter~\ref{ch:budget-relay}, or extend to the branching relays of
Chapter~\ref{ch:fanout}, whose fan-out structure introduces the further
complication of \cref{prop:silent-disagreement} and is not treated by the
linear construction given here. Appendix~\ref{app:open-questions} is
updated to reflect exactly this scope.
